
Morse theory allows us to extract from a manifold information about its topological invariants by studying the critical points of real-valued differentiable functions.

Given a differentiable function $f:M\to \mathbb{R}$ over a manifold $M$, Morse's lemma establishes that around every non-degenerate critical point there exists a local coordinate system with respect to which $f$ is expressed as a quadratic form. This, in turn, implies that non degenerate critical points are isolated and constitutes the foundation of Morse Theory. It is said that a function of this form is a Morse function if all its critical points are non degenerate.

If all sublevel sets of a Morse function $f$ are compact, it is possible to relate the homotopy type of a manifold with that of a CW-complex with as many cells of dimension $\lambda$ as critical points of index $\lambda$ $f$ has. In addition, Morse inequalities determine bounds for the Betti numbers of a manifold and let us compute its Euler-Poincaré characteristic relying only on the number of critical points of each possible index of a Morse function.

In this thesis we will develop the foundations of this theory following primarily the book ``Morse Theory'' by John Milnor \cite{milnor1963}.